% -*- mode: latex -*- %

\newcommand{\lipobj}{M}

\section{\ppipptitle~for general M-estimators}
\label{sec:ppipp-m-estimation}





Our results on generalized linear models follow from more general
results on M-estimators, which we turn to here. We consider convex
losses $\loss_\theta(x, y)$, and we continue to write
\begin{equation*}
  \theta\opt \defeq \argmin_\theta \E [\loss_\theta(X,Y)]
\end{equation*}
as the target of interest.
%
Our main theorem applies to what we term \emph{smooth enough} losses, which
roughly satisfy that $\loss_\theta$ is differentiable at $\theta\opt$ with
probability 1 and is locally Lipschitz; see
Definition~\ref{definition:smooth-enough} in
Appendix~\ref{sec:proof-normality}.
%
%% We use the short-hand notation $\ell_{\theta,i} =
%% \ell_{\theta}(X_i,Y_i)$, $\ell_{\theta,i}^f = \ell_{\theta}(X_i,f(X_i))$,
%% and $\widetilde \ell_{\theta,i}^f = \ell_{\theta}(\widetilde
%% X_i,f(\widetilde X_i))$. We will also sometimes use $\ell_{\theta}$ and
%% $\ell_{\theta}^f$ to denote the respective losses applied to a generic
%% draw from $\P$, e.g. $\ell_{\theta} \equiv \ell_{\theta,1}$ and
%% $\ell_{\theta}^f \equiv \ell^f_{\theta,1}$. We use analogous notation for
%% gradients, $\nabla\ell_{\theta,i}, \nabla\ell_{\theta,i}^f,
%% \nabla\widetilde \ell_{\theta,i}^f$, etc.
%
Smooth enough losses admit the following general asymptotic normality
result, whose proof we defer to Appendix~\ref{sec:proof-normality}. We use
the short-hand notation $\nabla\ell_{\theta} = \nabla\ell_{\theta}(X,Y)$ and
$\nabla \ell_{\theta}^f = \nabla \ell_{\theta}(X,f(X))$.

\begin{theorem}
  \label{theorem:normality}
  Assume that $\hat\lambda = \lambda + o_P(1)$ and that
  $\frac{n}{N}\rightarrow r$ for some $r\geq 0$. Let the losses
  $\loss_\theta$ be
  smooth enough (Definition~\ref{definition:smooth-enough}) and let $H_{\theta\opt} \defeq \nabla^2 L(\theta\opt)$.
  Define the covariance matrices
  \begin{align*}
    V_{f,\theta\opt}^\lambda \defeq \lambda^2 \cov(\nabla \loss_{\theta\opt}^f)\text{ and } V_{\Delta,\theta\opt}^\lambda  \defeq \cov( \nabla \loss_{\theta\opt} -  \lambda \nabla \loss_{\theta\opt}^f ).
  \end{align*}
  If the
  consistency $\thetaPP_{\hat \lambda} \stackrel{p}{\to} \theta\opt$ holds,
  then
  \begin{equation*}
    \sqrt{n}(\thetaPP_{\hat \lambda} - \theta\opt) \cd \normal\left(0, \Sigma^\lambda\right),
  \end{equation*}
  where 
  \begin{equation}
  \label{eqn:sigma-lambda}	
\Sigma^\lambda \defeq H_{\theta\opt}^{-1} \left(r \cdot
    V_{f,\theta\opt}^\lambda +
    V_{\Delta,\theta\opt}^\lambda
    \right) H_{\theta\opt}^{-1}.
    \end{equation}
\end{theorem}

The conditions sufficient for Theorem \ref{theorem:normality} to hold are
standard and fairly mild. As is familiar from classical asymptotic theory,
the one that is least straightforward to verify is the consistency of
$\thetaPP_{\hat \lambda}$, so here we give a few sufficient conditions. At a
high level, if $\LPP_\lambda$ is convex---as in the case of mean estimation
or GLMs---or the domain $\Theta$ of $\theta$ is compact, consistency holds
under mild additional regularity. Below we state the guarantee of
consistency for convex $\LPP_\lambda$ (providing a proof in
Appendix~\ref{sec:proof-consistency-convex}); for the treatment when
$\Theta$ is compact, see \citet[Chapter 5]{VanDerVaart98}.

\begin{proposition}
  \label{prop:consistency-convex}
  Assume that $\LPP_\lambda(\theta)$ is convex in $\theta$ with probability
  $1$ and that $\theta\opt$ is a unique minimum of $L(\theta)$.  Suppose
  also that $\loss_{\theta}$ is locally Lipschitz near $\theta\opt$
  (in the sense of
  Definition~\ref{definition:smooth-enough}\ref{item:locally-lipschitz}).
  Then, $\thetaPP_{\lambda} \cp
  \theta\opt$. If $\hat{\lambda} \cp \lambda$, then
  so long as $\P(\LPP_{\hat{\lambda}} ~\mbox{is~convex}) \to 1$, the same
  result holds.
\end{proposition}
\noindent
Because GLMs~\eqref{eq:glm-loss} and~\eqref{eq:exponential-family-loss}
have a convex loss $\LPP_\lambda(\theta)$
for $\lambda \in [0, 1]$, Theorem~\ref{theorem:normality} and
Proposition~\ref{prop:consistency-convex} combine to imply the asymptotic
normality result of Corollary~\ref{cor:normality_glms}.

With Theorem \ref{theorem:normality} in hand, we can generalize \ppipp developed for GLMs to general convex
M-estimators. In Algorithm~\ref{alg:general_ppi} we state the main
procedure, and Corollary~\ref{cor:validity_general} gives its validity.
Analogously to last section, $\widehat\Cov_n(\nabla
\loss_{\theta}),\widehat\Cov_{N+n}(\nabla \loss^f_{\theta})$, and
$\widehat\Cov_n(\nabla \loss_{\theta} - \hat \lambda \nabla
\loss_{\theta}^f)$ denote empirical covariance matrices computed using the
labeled data or the labeled and unlabeled data combined.


\begin{algorithm}[t]
\setstretch{1.3}
\caption{\ppipp: Prediction-powered inference for general M-estimators}
\label{alg:general_ppi}
\begin{algorithmic}[1]
  \Require labeled data $\{(X_i, Y_i)\}_{i=1}^n$, unlabeled data
  $\{\Xt_i\}_{i=1}^N$, predictive model $f$, error level $\alpha\in~(0,1)$,
  coefficient $j\in[d]$ of interest
  \State Select tuning parameter $\hat\lambda$
  \Comment{set tuning parameter}
  \State $\thetaPP_{\hat\lambda} = \argmin_\theta \LPP_{\hat\lambda}(\theta)$
  \Comment{prediction-powered point estimate}
  \State $\widehat H = \frac{1}{n} \sum_{i=1}^n \nabla^2 \ell_{\thetaPP_{\hat\lambda}}(X_i, Y_i)$
  \State $\widehat V_f = \hat \lambda^2 \widehat\Cov_{N+n} (\nabla \ell_{\thetaPP_{\hat\lambda}}^f)$ 
  \State $\widehat V_\Delta = \widehat\Cov_{n} (\nabla \ell_{\thetaPP_{\hat\lambda}} - \hat\lambda \nabla \ell_{\thetaPP_{\hat\lambda}}^f)$ 
  \State $\widehat\Sigma = \widehat H^{-1}( \frac{n}{N} \widehat V_f + \widehat V_\Delta) \widehat H^{-1}$ \Comment{covariance estimate}
  \Ensure 
  prediction-powered confidence set $\CPP_\alpha = \left(\thetaPP_{\hat\lambda, j} \pm z_{1-\alpha/2} \sqrt{\frac{\widehat\Sigma_{jj}}{n}} \right)$
\end{algorithmic}
\end{algorithm}


\begin{corollary}
  \label{cor:validity_general}
  Let the assumptions of Theorem~\ref{theorem:normality} hold and
  assume $\nabla^2 \loss_\theta$ is continuous in $\theta$.
  Then the
  prediction-powered confidence interval in Algorithm \ref{alg:general_ppi}
  has valid coverage: for any $j \in [d]$,
  \begin{equation*}
    \lim_{n} \P\left(\theta\opt_j \in \CPP_\alpha \right) = 1-\alpha.
  \end{equation*}
\end{corollary}

If we seek a confidence set for the whole vector $\theta\opt$, rather than a
single coordinate, many choices are
possible for a resulting confidence set. The ellipse
$\mc{E}_{\Sigma,\alpha} = \{v \in \R^d \mid v^\top \Sigma^{-1} v \le
\chi^2_{d,1-\alpha}\}$, where $\chi^2_{d,1-\alpha}$ denotes the $(1-\alpha)$-quantile
of the $\chi^2$ distribution with $d$ degrees of freedom,
satisfies $\P(Z \in \mc{E}_{\Sigma,\alpha}) = 1 - \alpha$ for
$Z \sim \normal(0, \Sigma)$. Thus
we can modify Algorithm~\ref{alg:general_ppi} to return
\begin{equation*}
  \CPP_\alpha = \thetaPP_{\hat\lambda}
  + \mc{E}_{\what{\Sigma} / n, \alpha},
\end{equation*}
which yields $\lim_n \P(\theta\opt \in \CPP_\alpha) = 1 - \alpha$.
Alternatively, we can apply the union bound, which gives rectangular sets of the form
\begin{equation*}
  \CPP_\alpha = \left[\thetaPP_{\hat\lambda,j} \pm \frac{z_{1-\alpha/(2d)}}{\sqrt{n}}
    \sqrt{\what{\Sigma}_{jj}} \right]_{j = 1}^d.
\end{equation*}
These sets guarantee $\liminf_n \P(\theta\opt \in \CPP_\alpha) \ge 1 - \alpha$.

